\documentclass[prl,aps,reprint,twocolumn,10pt,superscriptaddress,floatfix]{revtex4-2}
\usepackage[LGR,T1]{fontenc}
\usepackage{amsmath, amssymb, amsthm, mathtools, bbm, physics}
\usepackage{dsfont}
\usepackage[dvipsnames]{xcolor}
\usepackage[colorlinks=true, pdfborder={0 0 0}]{hyperref}
\definecolor{googleblue}{RGB}{34, 0, 204}
\definecolor{panblue}{RGB}{0,24,150}
\definecolor{carmine}{RGB}{150, 0, 24}
\hypersetup{
unicode=true,
colorlinks=true,
linkcolor=carmine,
citecolor=googleblue,
urlcolor=panblue,
anchorcolor=OliveGreen}

\newcommand{\tmp}[1]{\textcolor{red}{#1}}
\newcommand{\pedro}[1]{{\color{purple} #1}}
\newcommand{\bereket}[1]{{\color{blue} #1}}
\newcommand{\elie}[1]{{\color{RawSienna} #1}}

\usepackage{newtxtext}

\begin{document}

\title{\tmp{[Title: What does topological field theory require of probabilistic theories?]}}

\author{Pedro Lauand}
\affiliation{Perimeter Institute for Theoretical Physics, Waterloo, Ontario, Canada}

\begin{abstract}
\tmp{[Abstract.]}
\end{abstract}

\maketitle

\section{Introduction}

% Narrative map: ../syntheses/narratives/intro-narrative.md (B1--B9). Literature: ../syntheses/literature-map.md.
% Drafting beat by beat with Pedro; approved beats get uncommented here.

% B1+B2 (one paragraph): operational stance; GPTs/OPTs as the framework.
% (Approved 2026-08-31.)
\textcolor{blue}{Quantum theory is one of the most empirically successful
theories in science. It describes its experiments probabilistically: the
theory assigns probabilities to the possible outcomes of the measurements one
can perform upon physical systems. Every empirical claim of the theory
therefore is described by collecting measurement statistics of repeated
experimental runs. Despite its success, the foundational aspects of its
mathematical formulation have remained elusive~\tmp{[cite]}. In its textbook
presentation, quantum theory (QT) is expressed through complex vectors in a
Hilbert space and self-adjoint operators---an abstract structure whose
postulates arguably carry no direct physical motivation. Over the last
decades, however, this description has been reformulated into an
\emph{operational view}, i.e.\ to consider as its primitive notions
``self-evident'' features of general laboratory situations and derive the
abstract formalism from simple premises over these primitive notions. This
shift to the operational perspective allows one not to refer to the specific
mathematical framework of QT, but rather understand QT as a consequence of
natural information-theoretic constraints over the mathematical framework of
Generalized Probabilistic Theories (GPTs)~\cite{Barrett2007, Muller2021,
Plavala2023} in which QT can be situated as one theory among many
others~\cite{Hardy2001, CDP2011, MasanesMuller2011}.}

% B2b: reconstructions; SR as role model. (Approved 2026-08-31; grammar-only
% fixes: "also known as" clause moved next to its noun.)
\textcolor{blue}{Based on these ideas, the next wave of axiomatizations of QT,
also known as \emph{reconstructions} of QT, started. The most famous
historical example that is often cited as a role model for reconstructions of
QT is special relativity. In this example, Einstein showed that the Lorentz transformations
can be derived from essentially two simple physical principles: the constancy
of the speed of light and the relativity principle. This was an important
step towards improving our conceptual understanding of the Lorentz
transformations. Reconstructions of QT achieve a similar goal: by considering
information-processing assumptions, such reconstructions try to make the
abstract formalism of QT more physically motivated. The first reconstruction
of this flavor was the seminal work of Hardy~\cite{Hardy2001}, which later
inspired new solutions to the reconstruction
problem~\cite{CDP2011, MasanesMuller2011}. Each of these reconstructions
selects quantum theory out of the landscape of GPTs as the unique GPT that
satisfies a set of simple principles.}

% B3 (opening): NS as the most famous principle. (Approved 2026-09-01.)
\textcolor{blue}{The most famous physical principle in GPTs is the
\emph{no-signalling} principle~\cite{PopescuRohrlich1994}, which can be
understood as the weakest condition any ``reasonable'' physical theory must
satisfy, in the sense of being compatible with special
relativity~\cite{Gisin2020}. The no-signalling principle can be formulated as
the requirement that operations on separate systems
commute~\cite{Barrett2007}. This condition guarantees that when two devices
operate at spacelike separated locations, their correlations respect
relativistic causality. In this sense, the no-signalling principle can be
understood as the compatibility of relativistic causality with the operational
theory.}

% 2026-09-11 (agent, on Pedro's instruction): B3's closing question retired and
% replaced by the sharper question now stated in the next block. Old wording,
% approved 2026-09-01, kept here for reversion:
% This reasoning then suggests a natural question: are there other physical
% principles that spacetime compatibility requires from any probabilistic
% theory?

% New beginning (2026-09-01): NS = kinematical picture (points at spacelike
% separation) vs locality in TFT. Draft, awaiting approval.
% 2026-09-11 (agent, on Pedro's instruction): the programme's question added
% here, in Pedro's own words from the session of 2026-09-11 ("what does a
% probabilistic theory require to be distributed over space?"), replacing B3's
% closing question. Drafted, awaiting approval. "Space" rather than
% "spacetime" is deliberate (Pedro, 2026-09-11): distributed over space, then
% evolving in time.
\textcolor{blue}{The no-signalling principle considers the kinematical
picture of relativity: systems are embedded as points of spacetime, and
points held at spacelike separation cannot instantaneously influence each
other's measurement statistics. The only feature of spacetime this invokes
is the causal relation between a pair of points. This suggests a question
more precise than the compatibility of a theory with relativistic causality:
what does a probabilistic theory require in order to be distributed over
space? Another crucial feature of spacetime physics is the
\emph{locality of action}~\cite{Einstein1948}:
between points that \emph{are} causally connected, influences propagate
locally through spacetime, by a dynamical law.}

% RETIRED (kept for salvage): earlier B4 pivot — "modern view of spacetime
% physics", RG genericity, AQFT premise (i)/(ii) logic. Superseded by the
% new beginning above; RG/AQFT material may return in B7/B8.
% \textcolor{blue}{Here, we shed light on this question by incorporating a
% modern view of spacetime physics into the GPT framework. The no-signalling
% principle borrows from special relativity only its causal structure:
% spacetime enters as a record of which events may influence which, and
% constrains the correlations of devices held at spacelike separation. Modern
% spacetime physics, however, goes beyond this kinematic picture. Its degrees
% of freedom are \emph{fields}---physical quantities distributed over spacetime
% itself and governed by a local dynamical law~\cite{Einstein1948}. Moreover,
% this description is generic rather than fundamental: provided that discrete
% systems appear continuous at long distances, the renormalization group
% implies that microscopic details wash out and low-energy physics flows to a
% continuum field theory~\cite{WilsonKogut1974, Wilson1975, WeinbergEFT,
% Polchinski1992}.}
%
% \textcolor{blue}{This shift changes what compatibility with spacetime demands
% of a probabilistic theory. Once the degrees of freedom of a GPT are
% distributed over spacetime, two independent premises arise, mirroring the
% causal axioms of axiomatic quantum field
% theory~\cite{StreaterWightman, HaagSchroer1962, FewsterRejzner2019}:
% (i) degrees of freedom on spacelike separated regions are independent, so
% that systems compose in \emph{parallel}; and (ii) a local dynamical law
% propagates the degrees of freedom on a region to its causal development, so
% that systems compose in \emph{sequence}. The no-signalling principle
% implements premise (i) alone---indeed, the GPT tensor product is derived
% precisely from the commutation of operations on separate
% systems~\cite{Barrett2007}. Premise (ii), the locality of dynamics, has never
% been imported into the GPT framework; it is the premise from which our
% principles follow.}

% B6: reconceptualization -- continuity/RG justifies FT description.
\tmp{[B6]}

\tmp{[Warning]}

% B7: low-energy/gapped limit justifies TFT; Atiyah--Segal skeleton.
\tmp{[B7]}

% B8: the new question, mirroring NS.
\tmp{[B8]}

% B9: results.
\tmp{[Declare results.]}

\bibliography{draft}

\end{document}
